\documentclass[10pt,UTF8]{article}


\usepackage{geometry}

\usepackage[scheme=plain]{ctex}

\geometry{a4paper,scale=0.9}

\usepackage{bm} % 数学粗体（避免斜体被取消）

% \usepackage{makecell}
% \usepackage{multirow} % 表格多行合并

\usepackage{booktabs} % 表格美化

% \usepackage{graphicx} %插入图片的宏包
% \usepackage{subfigure}
% \usepackage{float} %设置图片浮动位置的宏包

\usepackage{amsmath}
\usepackage{esint} % 提供更丰富的积分符号
\usepackage{amssymb}
\usepackage{mathtools}

\usepackage{physics}

% 选择数学字体
% 符合 ISO 
\usepackage[utopia]{mathdesign}
\usepackage{isomath}

\renewcommand{\vec}[1]{\vectorsym{#1}}
\newcommand{\mat}[1]{\matrixsym{#1}}

\usepackage{fontspec} % 提供字体选择接口
\setmainfont{Times New Roman} % 设置正文字体为 Times New Roman

% \usepackage{siunitx}
% % 关闭警告
% \ExplSyntaxOn
% \msg_redirect_name:nnn { siunitx } { physics-pkg } { none }
% \ExplSyntaxOff

% \usepackage{bussproofs} % 自然演绎推理树

% 交叉引用
% 关闭边框，设置颜色
\usepackage[colorlinks=true, linkcolor=darkgray, citecolor=teal,
urlcolor=blue]{hyperref}
% 使 \cref 和 \Cref 可用
\usepackage{cleveref}

\usepackage{bookmark}  % 生成 PDF 大纲

\allowdisplaybreaks[3]  % 允许换页


%%%%%%%%%%%%%%%%%%%%%%%% tcolorbox %%%%%%%%%%%%%%%%%%%%%%%%%%
\usepackage[most]{tcolorbox}

\newenvironment{tipBox}
{
  \begin{tcolorbox}[colback=green!5!white,colframe=green!75!black,parbox=false,before
  upper=\par,before lower=\par]}
  {
\end{tcolorbox}}

\newenvironment{conceptBox}
{
\begin{tcolorbox}[colback=blue!5!white,colframe=blue!75!black]}
  {
\end{tcolorbox}}

\newenvironment{theoremBoxNoIndent}
{
\begin{tcolorbox}[colback=yellow!5!white,colframe=yellow!75!black]}
  {
\end{tcolorbox}}

\newenvironment{definitionBox}
{
  \begin{tcolorbox}[colback=blue!5!white,colframe=blue!75!black,parbox=false,before
  upper=\par,before lower=\par]}
  {
\end{tcolorbox}}

\newenvironment{theoremBox}
{
  \begin{tcolorbox}[colback=yellow!5!white,colframe=yellow!75!black,parbox=false,before
  upper=\par,before lower=\par]}
  {
\end{tcolorbox}}
%%%%%%%%%%%%%%%%%%%%%%%%%%%%%%%%%%%%%%%%%%%%%%%%%%%%%%%%%%%%

%%%%%%%%%%%%%%%%%%%%% amsthm %%%%%%%%%%%%%%%%
\usepackage{amsthm}

\theoremstyle{definition}
\newtheorem{definition inner}{Definition}[section]
\newenvironment{definition}[1][]
{
  \begin{definitionBox}
  \begin{definition inner}[#1]\hfill\par}
    {
    \end{definition inner}
\end{definitionBox}}

\theoremstyle{definition}
\newtheorem{theorem inner}{Theorem}[section]
\newenvironment{theorem}[1][]
{
  \begin{theoremBox}
  \begin{theorem inner}[#1]\hfill\par}
    {
    \end{theorem inner}
\end{theoremBox}}

\theoremstyle{definition}
\newtheorem*{proof inner}{\textit{Proof}}
\renewenvironment{proof}[1][]
{
\begin{proof inner}[#1]\hfill\par}
  {\qed\vphantom{.}
\end{proof inner}}

\theoremstyle{plain}
\newtheorem{tip inner}{Tip}[section]
\newenvironment{tip}[1][]
{
  \begin{tipBox}
  \begin{tip inner}[#1]\hfill\par}
    {
    \end{tip inner}
\end{tipBox}}

\theoremstyle{remark}
\newtheorem*{remark}{\textit{Remark}}
%%%%%%%%%%%%%%%%%%%%%%%%%%%%%%%%%%%%%%%%%%%%%%


\newcommand{\iu}{\mathrm{i}}  % 虚数单位 i


\begin{document}

\part*{Systems of First Order Linear ODEs}

\section{First order ODE systems}

The general system of first order equations involving
\(n\) dependent variables defined on an interval \(I \subseteq \mathbb{R}\)
can be written as
\begin{equation}
  \begin{gathered}
    \begin{cases}
      y_1'(t) = F_1(t, y_1, \ldots, y_n) \\
      y_2'(t) = F_2(t, y_1, \ldots, y_n) \\
      \qquad\vdots \\
      y_n'(t) = F_n(t, y_1, \ldots, y_n)
    \end{cases} \\
    \text{with initial conditions}\quad
    y_1(t_0) = x_1, \quad
    y_2(t_0) = x_2, \quad
    \ldots, \quad
    y_n(t_0) = x_n,
  \end{gathered}
  \qfor t\in I.
  \label{eq: system of 1st order ODEs}
\end{equation}
The solution to the system can be expressed in vector form as
\begin{equation*}
  \vec{y}(t) =
  \begin{pmatrix}
    y_1(t) \\
    y_2(t) \\
    \vdots \\
    y_n(t)
  \end{pmatrix}
\end{equation*}

\begin{theorem}[Existence and Uniqueness]
  {}\label{thm: Existence and Uniqueness for 1st order ODE Systems}
  For the system of first order ODEs~(\ref{eq: system of 1st order ODEs}),
  if the functions \(F_1, F_2, \ldots, F_n\) and the partial derivatives
  \(\displaystyle\pdv{F_i}{y_j}\) for all pairs \(i, j = 1, 2, \ldots, n\)
  are continuous on a region \(R\) defined by
  \begin{equation*}
    R := [t_0 - a, t_0 + a] \times [x_1 - b, x_1 + b] \times \cdots
    \times [x_n - b, x_n + b] \subseteq \mathbb{R}^{n+1},
  \end{equation*}
  where constants \(a, b > 0\), then there exists a unique solution
  \(\vec{y}(t)\) for \(t \in (t_0 - h, t_0 + h)\) to the IVP,
  where
  \begin{equation*}
    h = \min \{a, b/M\} \qand
    M = \max\limits_{\vec{v} \in R}\{
      \abs{F_1(\vec{v})}, \abs{F_2(\vec{v})}, \ldots, \abs{F_n(\vec{v})}
    \}.
  \end{equation*}
\end{theorem}

Any \(n\)-th order ODE can be rewritten as a system of \(n\) first order ODEs.
Given an \(n\)-th order ODE
\begin{equation*}
  y^{(n)} = F(t, y, y', \ldots, y^{(n-1)}),
\end{equation*}
setting \(z_1 = y, z_2 = y', \ldots, z_n = y^{(n-1)}\), we obtain the system
\begin{equation*}
  \begin{cases}
    z_1' = z_2 \\
    z_2' = z_3 \\
    \qquad\vdots \\
    z_{n-1}' = z_n \\
    z_n' = F(t, z_1, z_2, \ldots, z_n)
  \end{cases}
\end{equation*}
which is a system of \(n\) first order ODEs.

\begin{definition}[First-Order Linear System]
  {}\label{def: 1st-order linear system}
  A system of first order ODEs is called \textbf{linear} if it can be
  expressed in the form
  \begin{equation}
    \vec{y}'(t) = \mat{P}(t) \vec{y}(t) + \vec{g}(t) \quad \qfor t\in I,
    \label{eq: 1st-order linear system}
  \end{equation}
  where
  \begin{equation*}
    \vec{y}(t) =
    \begin{pmatrix}
      y_1(t) \\
      y_2(t) \\
      \vdots \\
      y_n(t)
    \end{pmatrix}
    \qquad
    \vec{g}(t) =
    \begin{pmatrix}
      g_1(t) \\
      g_2(t) \\
      \vdots \\
      g_n(t)
    \end{pmatrix}
    \qquad
    \mat{P}(t) =
    \begin{pmatrix}
      p_{11}(t) & p_{12}(t) & \cdots & p_{1n}(t) \\
      p_{21}(t) & p_{22}(t) & \cdots & p_{2n}(t) \\
      \vdots & \vdots & \ddots & \vdots \\
      p_{n1}(t) & p_{n2}(t) & \cdots & p_{nn}(t)
    \end{pmatrix}
  \end{equation*}
\end{definition}

\begin{definition}[Homogeneous First-Order Linear System]
  {}\label{def: Homogeneous 1st-order linear system}
  A first order linear system is called \textbf{homogeneous} if
  \(\vec{g}(t) = \vec{0}\) for all \(t\in I\), i.e., it can be expressed as
  \begin{equation*}
    \vec{y}'(t) = \mat{P}(t) \vec{y}(t) \quad \qfor t\in I.
  \end{equation*}
\end{definition}

\begin{theorem}[Existence and Uniqueness for First-Order Linear Systems]
  {}\label{thm: Existence and Uniqueness for 1st-order linear systems}
  For the first-order linear system~(\ref{eq: 1st-order linear system}),
  if all entries of \(\mat{P}(t)\) and \(\vec{g}(t)\) are continuous on
  an interval \(I\), then for any \(t_0 \in I\) and \(\vec{y}_0 \in
  \mathbb{R}^n\),
  there exists a unique solution \(\vec{y}(t)\) to the IVP
  \begin{equation*}
    \begin{cases}
      \vec{y}'(t) = \mat{P}(t) \vec{y}(t) + \vec{g}(t) \\
      \vec{y}(t_0) = \vec{y}_0
    \end{cases}
    \quad \qfor t\in I.
  \end{equation*}
\end{theorem}

\end{document}
